% main4-body.tex — Clean review article body.
% All internal file paths, script references, and process artifacts removed.

\section{Introduction}
\label{sec:intro}

An agricultural vision system may observe the same tree from several positions, frames, or sensors. The operational output is not the number of detections emitted by those images. It is an inventory of unique fruit instances, grouped by an attribute class. For a tree with instance set~$I$, the desired count for class~$c$ is the number of instances whose assigned attribute is~$c$. A detection that is missed, merged with a neighbour, assigned the wrong class, or counted again from another view produces a different error. These errors cannot be diagnosed from image-level accuracy alone.

This review addresses the task as a design problem:
\begin{quote}
\emph{How can many observations of one tree be converted into a class-wise inventory in which each fruit or fruit bunch is counted once?}
\end{quote}
The principal agricultural case is the oil-palm fresh fruit bunch (FFB). The class attribute is not restricted to maturity. Maturity is one example; size, quality, disease, defect, cultivar, or harvest readiness are also possible attributes. All of them remain instance-dependent because a class-wise count is undefined until an individual fruit or bunch has been localized and assigned an identity.

The scope is agricultural fruit perception, but the mechanism search is deliberately wider. Tracking, person or vehicle re-identification, multi-camera association, visual odometry, structure from motion (SfM), and 3D object detection were developed largely outside orchards. They are included as transferable evidence when they answer a mechanism question about identity, geometry, or duplicate rejection. Their results are not presented as agricultural performance.

The review makes four contributions. First, it reports a reproducible Scopus and OpenAlex search with query strings, date bounds, selection rules, and a deterministic deduplication record. Second, it organizes the task around six identity mechanisms and the acquisition assumptions each mechanism requires. Third, it positions previous oil-palm, fruit-counting, and fruit-ripeness reviews against the cross-view identity problem. Fourth, it provides a row-level synthesis matrix in which method, modality, evaluation context, identity mechanism, class attribute, finding, and transfer boundary are visible together.

The paper does not propose one architecture as the universal solution. It asks when a simple statistical correction, temporal tracking, appearance matching, geometric reconstruction, or learned multi-view model is justified by the acquisition regime. This distinction is important for oil palm because the same tree can be observed from four or eight positions while sensor calibration, depth quality, illumination, and occlusion vary at the same time.


\section{Review methodology}
\label{sec:protocol}

\subsection{Review type and questions}

This work is a structured design-space review reported with a PRISMA-style search and selection trace. It is not a meta-analysis. The included tasks, object definitions, sensor modalities, and evaluation metrics are too heterogeneous for a pooled effect size. The protocol was not registered before execution. The review questions are:

\begin{enumerate}[leftmargin=*,itemsep=1pt]
  \item What operations are required to turn repeated observations into a unique, class-wise inventory?
  \item Which identity mechanisms have been used for fruit or for transferable non-agricultural objects?
  \item Which assumptions about overlap, motion, calibration, appearance, and view order make each mechanism plausible?
  \item What evidence exists for oil-palm and other tree-fruit systems, and what remains a transfer hypothesis?
  \item Which measurements would distinguish a better detector from a better inventory system?
\end{enumerate}

\subsection{Search strategy}

Scopus was used as the subscription index available in the institutional session. OpenAlex was used as an open complementary index. Web of Science was not available in the recorded run; this deviation is stated rather than hidden. Google Scholar was reserved for possible backward or forward snowballing and was not used as a primary count source because its result export is not stable. The date window was 2015--2026.

Seven query families were used, each combining topic terms with a subject-area filter (Agriculture, Computer Science, Engineering, and Multidisciplinary). Their targets are: Q1, multi-view fruit inventory; Q2, cross-view identity and duplicate rejection; Q3, plant geometry and multimodal sensing; Q4, per-instance class attributes; Q5, previous reviews; Q6, oil-palm vision and machine learning; Q7, single-view fruit counting and yield estimation. The exact search strings are reproduced in Appendix~\ref{app:queries}.

\begin{table}[t]
\centering
\caption{Search yields before cross-source deduplication.}
\label{tab:searchcounts}
\small
\begin{tabular}{lrrl}
\toprule
Query & Scopus & OpenAlex & Main focus\\
\midrule
Q1 & 2{,}001 & 1{,}849 & Multi-view fruit inventory\\
Q2 & 2{,}789 & 1{,}991 & Identity and duplicates\\
Q3 & 5{,}648 & 6{,}422 & Plant geometry and sensors\\
Q4 & 2{,}422 & 2{,}757 & Per-instance attributes\\
Q5 & 662 & 1{,}143 & Previous reviews\\
Q6 & 995 & 1{,}177 & Oil-palm vision\\
Q7 & 1{,}205 & 1{,}317 & Single-view fruit load\\
\midrule
Total & 15{,}722 & 16{,}656 & \\
\bottomrule
\end{tabular}
\end{table}

Known-item checks confirmed that the subject-area filter did not remove three anchor studies: SawitMVC in Q1, Gené-Molá et al. in Q3, and MangoYOLO in Q7. These checks demonstrate filter safety but not index completeness.

\subsection{Eligibility criteria}

A source was eligible for the focused evidence register when all five inclusion conditions were met:

\begin{description}[leftmargin=*,style=nextline,itemsep=2pt]
  \item[IC1.] Peer-reviewed research article, conference paper, or openly available preprint with a verifiable full text and a quantitative result.
  \item[IC2.] Visual output at the per-instance level: object detection, instance segmentation, tracking, localization, or cross-observation association.
  \item[IC3.] Publication year 2015--2026, except for foundational citations used only for context.
  \item[IC4.] Written in English or Indonesian.
  \item[IC5.] Full text identifies a dataset or acquisition protocol and an evaluation metric or quantitative comparison.
\end{description}

The exclusion criteria were: EC1, output only at image or scene level without per-instance identity; EC2, no quantitative evaluation; EC3, full text unavailable; EC4, duplicate or version of an already retained record; EC5, editorial, abstract-only, poster, patent, or non-research document; EC6, outside the date range.

\subsection{Screening and selection}

The raw exports were kept unchanged. Within each query, exact duplicate identifiers were removed. Across sources, the primary deduplication key was a normalized DOI; records without a DOI used normalized title plus year. A title-year collision was not merged when author or venue evidence conflicted. The seven Scopus exports contained 15,722 rows and the seven OpenAlex exports contained 16,656 rows. Together they produced 32,378 raw rows, 32,222 after within-query deduplication, and 22,269 records after exact-key union. A conservative resolver processed 1,030 merge components and reduced the working master to 21,066 records. The complete deduplication decisions and screening records are available as supplementary material.

Title screening removed 243 records (242~EC5, 1~EC6), leaving 20,823 for abstract screening. Abstract screening removed a further 788 records (640~EC1, 148~EC5), leaving 20,035 candidates. Fig.~\ref{fig:flow} summarizes the selection trace.

\begin{figure}[t]
\centering
\small
\begin{tabular}{c}
\fbox{\parbox{0.78\columnwidth}{\centering 32{,}378 raw records\\Scopus 15{,}722 + OpenAlex 16{,}656}}\\[2pt]
$\downarrow$\\[-1pt]
\fbox{\parbox{0.78\columnwidth}{\centering 22{,}269 exact-key master records}}\\[2pt]
$\downarrow$\\[-1pt]
\fbox{\parbox{0.78\columnwidth}{\centering 21{,}066 resolved working-master records}}\\[2pt]
$\downarrow$\\[-1pt]
\fbox{\parbox{0.78\columnwidth}{\centering 20{,}035 abstract-screened candidates}}\\[2pt]
$\downarrow$\\[-1pt]
\fbox{\parbox{0.78\columnwidth}{\centering 250 shortlisted\quad$\rightarrow$\quad 44 included}}
\end{tabular}
\caption{Search and selection trace. The 20,035 candidates were ranked and shortlisted before targeted full-text review; this count represents the search pool, not a promise of individual paper reviews.}
\label{fig:flow}
\end{figure}

The 20,035 candidates were not treated as a mandatory reading list. A deterministic title-abstract scoring function ranked signals relevant to identity, inventory, duplicate resolution, re-identification, cross-view association, tracking, SfM, RGB-D, depth, fruit, FFB, and oil palm. Penalties reduced the priority of global yield or biomass outputs, canopy-only remote sensing, non-fruit targets, and image-level outputs. To prevent one model family from dominating, selection used diversified buckets for core identity, direct oil-palm evidence, fruit multi-view or 3D evidence, fruit-instance baselines, transfer mechanisms, and prior reviews. The resulting shortlist contained 250~records with a first targeted wave of 60. By the manuscript cutoff, 44~records had complete full-text evidence and 16 had been excluded after review.

\subsection{Evidence registers}

The manuscript keeps two registers separate. Register~A is the full Scopus and OpenAlex search, including all candidate and screening states. Register~B is the focused evidence register: studies whose full text was retrieved and whose dataset, protocol, quantitative evaluation, and mechanism fields were verified. The ranking and shortlist are selection aids inside Register~A; they are not inclusion decisions. This separation prevents the 20,035 search-pool records from being presented as 20,035 included studies.


\section{The inventory task as a design space}
\label{sec:designspace}

\subsection{From detections to inventory}

The task has five distinct operations. First, a detector or segmenter proposes an instance in each observation. Second, a class head attaches an attribute to that instance. Third, observations are associated when they refer to the same physical object. Fourth, duplicate observations are rejected or merged. Fifth, the remaining identities are aggregated into counts per class. The first two operations can be evaluated with precision, recall, average precision, and class confusion. The last three require identity and count metrics.

This decomposition changes the interpretation of a strong detector. A detector with high average precision on independent images may still overcount when a tree is filmed from both sides. A tracker may preserve identity through a short video but fail when the camera returns after a gap. A 3D association method may remove duplicates but become unusable when depth is uncalibrated. The design choice is therefore conditional on the acquisition regime.

\subsection{Six identity mechanisms}

Table~\ref{tab:mechanisms} defines the six mechanisms used in the synthesis.

\begin{table}[t]
\centering
\caption{Identity mechanisms and their minimum assumptions.}
\label{tab:mechanisms}
\small
\begin{tabularx}{\columnwidth}{@{}lY Y@{}}
\toprule
Code & Mechanism & Main assumption and risk\\
\midrule
M0 & Intra-view detection or segmentation & One view is sufficient. Hidden objects and repeated views are not resolved.\\
M1 & Statistical correction or fixed census & Duplicate structure is stable enough for a correction factor. The correction may not transfer across trees.\\
M2 & Appearance or re-identification & The instance retains discriminative texture, colour, or shape despite illumination and occlusion.\\
M3 & Temporal tracking & Consecutive observations overlap in time and motion is sufficiently smooth. Track fragmentation and drift remain possible.\\
M4 & Geometric or 3D association & Calibration, scale, or reliable reconstruction is available. Sensor failure and SfM cost are risks.\\
M5 & Learned multi-view association & Linked training data represent view, class, and occlusion variation. The model may overfit one acquisition pattern.\\
\bottomrule
\end{tabularx}
\end{table}

M0 is the no-association baseline. M1 uses a statistical correction or a fixed counting protocol without assigning a persistent identity. M2 matches local appearance or a learned re-identification embedding. M3 uses temporal continuity, motion, and track management. M4 uses calibrated geometry, depth, SfM, or a 3D map. M5 learns a multi-view association function directly from linked observations. A system can combine mechanisms, but its reported result must identify which mechanism is responsible for duplicate rejection.

\subsection{Acquisition assumptions}

The second design axis is the assumption set. A1~is the amount of overlap and occlusion. A2~is the camera baseline or view separation. A3~is motion smoothness or scene stationarity. A4~is calibration and metric scale. A5~is appearance distinctiveness. A6~is regularity of duplicate exposure, for example a fixed four-side census. A7~is the availability of identity supervision or an oracle association. M1~mainly depends on A6. M2~depends on A5. M3~depends on A3. M4~depends on A2 and A4. M5~depends on A7 and on coverage of the training acquisition pattern.

This two-axis view gives a precise meaning to the phrase ``multi-view deduplication.'' It is not a mandatory module. It is a mechanism whose value should be tested against the assumptions that survive in the field. A fixed four-side protocol may make M1 competitive when duplication is regular and metric depth is unreliable. An unstructured moving camera may require M3 or M4 because the number and order of observations are not fixed.

\subsection{Class attributes beyond maturity}

Maturity classification is common in FFB literature, but the identity dependency is more general. An instance may carry a maturity label, a size class, a defect label, a disease state, a cultivar label, or a harvest-readiness state. The inventory should first establish one identity record and then attach one or more attributes to that record. If the same bunch appears in three views with two maturity predictions, the system must report whether the conflict is resolved by confidence, temporal evidence, geometry, or manual adjudication. Summing per-view class predictions without this step is not a class-wise inventory.


\section{Evidence from agricultural and fruit systems}
\label{sec:direct}

\subsection{Oil-palm detection and classification}

The direct oil-palm literature is concentrated on detection, ripeness classification, dataset construction, and yield-related counting. Lai et al. reviewed FFB ripeness detection methods across sensing families, while Goh et al. reviewed non-destructive classification methods including computer vision, LiDAR, spectroscopy, thermal sensing, and laser-light backscatter~\cite{lai2023palmslr,goh2025palmreview}. These reviews establish that sensor choice and class definition are central problems. Their primary unit is ripeness assessment, not the persistent identity of one bunch across repeated views.

Naftali et al. compared YOLO and other detectors on an oil-palm plantation dataset and reported a depthwise-separable YOLOv8 variant as a real-time candidate~\cite{naftali2024palmcounter}. This is useful evidence for M0 detector selection, but it does not by itself measure whether a bunch detected in one frame is the same bunch detected in another frame. A detector benchmark and an inventory benchmark answer different questions.

Three dataset papers expose the acquisition variables needed for the design space. Goh et al. released 400 RGB images and 400 registered depth maps from four plantation locations, captured with an Intel RealSense D455f. The article reports a mean registration error of 1.8~cm at 3~m and 92.5\% expert inter-rater agreement for binary ripeness labels~\cite{goh2025palmdataset}. Suharjito et al. released plantation imagery captured from multiple angles with five categories: unripe, underripe, ripe, flower, and abnormal~\cite{suharjito2025palmdataset}. These datasets address sensor and class variation, but their published descriptions do not make a cross-view identity graph the primary evaluation object.

SawitMVC is directly aligned with the inventory problem. Its source article describes 3,992 smartphone images of 953 oil-palm trees from two plantations, captured from four or eight angular positions. The dataset uses four BBC maturity classes and a separate per-tree JSON layer with 9,823 unique bunches and cross-view identity links~\cite{indriani2026sawitmvc}. This makes a tree-level unique-count metric possible. It also exposes the central confounder: the number of visible bounding boxes is not the number of unique bunches.

Additional oil-palm studies in the evidence matrix include LiDAR-intensity scanning for ripeness~\cite{ledger_r_97003608b250cdc5}, UAV SfM for canopy height and per-palm inventories~\cite{ledger_m_084cfd5badbb63d3}, RGB-depth detection with YOLOv4~\cite{lai2022yolov4palm}, and several detector comparisons using ground-level images~\cite{ledger_r_19d96439ecdda18d,chen2022ffbdet,ledger_r_134385f2faed38a6}. Fluorescence imaging under controlled conditions also provides a sensor modality for attribute classification~\cite{ledger_r_1c81e223679e4bfe}. In every case the output is per-image detection or classification; none of these studies reports tree-level deduplication.

\subsection{Fruit counting and yield estimation}

Single-view fruit-load studies show the limit of a detector-only interpretation. Koirala et al. compared six existing architectures and developed MangoYOLO for mango detection, achieving an F1 of 0.968 and average precision of 0.983, while orchard load estimates still required a correction factor against hand-harvest counts~\cite{koirala2019deepfruit}. The correction is an M1 mechanism: it can be operationally useful without solving identity, but its validity depends on stable exposure and sampling assumptions.

Gené-Molá et al. made identity and 3D location explicit by combining Mask R-CNN instance segmentation with SfM, projecting image detections into a 3D point cloud, and removing false positives with a support vector machine. The study used 11~Fuji apple trees containing 1,455~apples. The F1 score increased from 0.816 for 2D detection to 0.881 for 3D detection and location, while the authors identified SfM processing time as a barrier to real-time use~\cite{genemola2020fruit3d}. This is direct evidence that geometric association (M4) can change the counting result.

Temporal tracking (M3) has been applied to orchard video. Zhang et al. combined deep-learning detection with a counting-region strategy for citrus, achieving a mean absolute error of 0.081 compared with 0.45 for conventional SORT~\cite{ledger_r_ee550a3e3a6fb318}. Gené-Molá et al. applied ByteTrack to apple video, reporting HOTA of 0.689 and yield-mapping MAPE of 7.47\%, while noting that appearance-based re-identification did not improve tracking for visually similar apples~\cite{ledger_r_2d972b94554c6cd1}. Multi-view 3D reconstruction using Semantic NeRF~\cite{ledger_m_1210e3ea00486b68} and RGB-D SLAM~\cite{ledger_r_b2e0e9aabad17a97} provide complementary geometric routes to consolidate repeated observations before counting.

A pear study combining LiDAR--camera fusion with FAST-LIO2 SLAM reported 96.2\% counting accuracy while retaining 53~double-counts linked to association threshold and SLAM drift~\cite{ledger_r_d4270ea8ab11f268}. Rapado-Rincón et al. evaluated multi-view tomato perception with HOTA, detection association, and count error, demonstrating that total count error alone can hide ID switches and false positives~\cite{ledger_r_700d34dae0f20d78}. Fu et al. used RGB and depth features with Faster R-CNN for apple detection in dense foliage, supporting the role of depth as an occlusion and localization cue~\cite{fu2020apple}.

\subsection{Positioning against previous reviews}

The reviews found by Q5 are adjacent rather than identical to this review. Table~\ref{tab:position} records the scope visible in the accessible article record and identifies whether cross-view identity is treated as the primary unit. The wording is deliberately bounded: ``not explicit'' means that the accessible scope and abstract do not make identity and duplicate rejection the central organizing question.

\begin{table}[t]
\centering
\caption{Positioning against previous reviews.}
\label{tab:position}
\scriptsize
\begin{tabularx}{\columnwidth}{@{}p{0.24\columnwidth}p{0.22\columnwidth}p{0.14\columnwidth}Y@{}}
\toprule
Review & Primary scope & Cross-view identity? & Position\\
\midrule
Lai et al.~\cite{lai2023palmslr} & FFB ripeness detection & Not explicit & Sensor context; does not organize counting mechanisms.\\
Goh et al.~\cite{goh2025palmreview} & FFB classification, non-destructive sensors & Not explicit & Modality limits; this review adds identity evaluation.\\
Naftali et al.~\cite{naftali2024palmcounter} & FFB detector comparison & Not explicit & Detector context; this review separates detection from inventory.\\
Naghipour et al.~\cite{naghipour2024fruitreview} & AI for fruit detection & Not explicit & Broad coverage; this review makes association a design axis.\\
Sapkota et al.~\cite{sapkota2024yoloagri} & YOLO in agriculture & Not explicit & Architecture taxonomy; this review is task-first.\\
Nordin and Misro~\cite{nordin2026palmslr} & Palm-oil maturity, deep learning & Not explicit & Maturity comparator; this review generalizes the class attribute.\\
\bottomrule
\end{tabularx}
\end{table}

The bounded conclusion is not ``no previous review exists.'' It is that the accessible review set is organized mainly by sensor, ripeness, detector family, or agricultural application. The cross-view identity of each physical fruit as the prerequisite for a class-wise inventory is not the primary comparison unit in that set. This is the gap addressed here.


\section{Transferable mechanisms for identity and duplicate rejection}
\label{sec:transfer}

\subsection{Appearance and temporal association}

M2 matches instances by appearance. A descriptor may use colour, texture, shape, or a learned embedding. Its strength is that it does not require metric depth. Its failure modes are predictable in fruit scenes: two bunches of the same maturity can look alike, illumination changes colour, and partial occlusion removes the most distinctive region. An evaluation must therefore report identity precision and recall under appearance changes, not only detector average precision. Weakly supervised cross-view embeddings from the vehicle re-identification domain show that identity labels alone can reduce viewpoint-induced variation~\cite{ledger_r_4b3d8d423dbd5d0e}, although the assumption that identity labels are available during training must be explicitly stated.

M3 uses time. A tracker predicts where an instance should appear in the next frame, associates detections, and terminates a track when evidence is insufficient. Video-based oil-palm ripeness work shows why temporal input is attractive for deployment~\cite{junior2023videopalm}. Temporal continuity is weaker when the camera moves around a tree and revisits a previously observed side. The report must state the maximum time gap and whether the camera path is ordered or arbitrary.

\subsection{Geometric and 3D association}

M4 associates observations through camera pose, depth, point clouds, or an SfM reconstruction. Gené-Molá et al. provide the clearest fruit example: a 2D instance mask is projected into a shared 3D representation so that repeated image detections can be consolidated~\cite{genemola2020fruit3d}. Park et al. describe a real-time orchard approach that incrementally builds 3D crop instances and uses 3D generalized intersection-over-union for association~\cite{park2026cropassociation}. The mechanism is directly aligned with duplicate rejection, but its field cost includes calibration, reconstruction stability, and sensor coverage.

Non-agricultural 3D detection studies provide transferable design components rather than fruit results. Camera and LiDAR fusion methods such as 3D-CVF and TransFusion show how multiple sensor views can be aligned at feature or decision level~\cite{yoo20203dcvf,bai2022transfusion}. Visual SLAM systems show how a moving camera can estimate pose and maintain a scene map, while dynamic-scene variants explicitly mask moving objects before mapping~\cite{bescos2018dynaslam}. These studies motivate a measurement question for agriculture: does a shared geometric frame reduce false duplicates after controlling for detector recall and sensor failure?

\subsection{Learned multi-view association}

M5 learns identity from linked observations. Training examples must state which detections belong to one physical object and which are different objects. Without this supervision, a model can learn a view-specific shortcut and still appear accurate on a random image split. SawitMVC's per-tree identity links create the type of supervision required for this experiment~\cite{indriani2026sawitmvc}. The necessary split is by physical tree or plantation block, not by randomly distributing frames from one tree across train and test.


\section{Modalities and fusion}
\label{sec:fusion}

RGB provides colour, texture, and fine appearance. Its weaknesses are shadows, specular highlights, camouflage, and scale changes. Depth provides a spatial cue, but ``depth'' is not one measurement. A calibrated sensor may provide metric range; stereo and SfM provide geometry relative to a reconstructed scene; a monocular network provides a learned or pseudo-depth ordering whose scale may not be metric. A manuscript must state which meaning is used.

Three fusion points are useful design categories. Early fusion concatenates RGB and depth at the input. It is simple but exposes the detector to missing or misregistered channels. Middle fusion extracts separate features and combines them at an intermediate layer, which permits modality-specific processing and quality gating. Late fusion combines detections, scores, or identity graphs after each modality has produced an estimate. RGB-depth object-detection studies demonstrate that the choice of fusion point is a measurable hypothesis rather than a universal rule~\cite{ophoff2019multimodal}. An oil-palm RGB-depth CNN and visual-SLAM system illustrates a direct agricultural attempt to combine appearance, depth, and localization~\cite{siong2021rgbdepthpalm}.

The design-space consequence is that sensor quality belongs in the evaluation, not only in the apparatus description. Each experiment should report the valid-depth ratio, RGB-depth registration error, missing-depth policy, and the fraction of detections supported by valid geometry. A fusion gain observed after dropping invalid frames is not comparable with a gain measured on all frames. The Goh dataset's reported registration error provides an example of the type of calibration evidence needed before interpreting a depth-based identity result~\cite{goh2025palmdataset}.


\section{Synthesis: from detector score to inventory evidence}
\label{sec:synthesis}

\subsection{Design rules}

The evidence supports five design rules:

\begin{enumerate}[leftmargin=*,itemsep=2pt]
  \item Report per-view detection separately from cross-view identity. A detector score cannot establish duplicate-free counting.
  \item Define the unit of truth at tree level. For each physical tree, store the unique instance count, per-view appearances, class labels, and visibility state.
  \item Treat the acquisition path as an experimental factor. A fixed four-side census, a continuous video, and an arbitrary multi-camera setup impose different assumptions.
  \item Test identity mechanisms against their required assumptions. M1 needs regular duplication; M3 needs temporal overlap; M4 needs geometry; M5 needs identity-linked training data.
  \item Report class-wise count error after identity resolution. A correct class label on a duplicate is still a wrong inventory.
\end{enumerate}

\subsection{Evaluation requirements}

The minimum evaluation bundle is: per-instance precision and recall, class-wise average precision, identity precision and recall, duplicate rate, missed-appearance rate, absolute and relative count error per tree, class-confusion error per unique instance, sensor-validity statistics, and latency. Where tracks are used, IDF1 or HOTA can complement track precision and recall. Where 3D association is used, calibration error and association distance must be reported. These metrics should be paired by physical tree so that a difficult canopy is not assigned to one method and an easy canopy to another.


\section{Research gaps and testable agenda}
\label{sec:agenda}

The first gap is an evaluation gap. Many studies report detection or ripeness classification, but a class-wise inventory requires identity links and a tree-level unique-count ground truth. SawitMVC demonstrates the data structure needed for this measurement, while the earlier fruit and FFB literature often reports image-level or frame-level results~\cite{indriani2026sawitmvc,naftali2024palmcounter,koirala2019deepfruit}.

The second gap is a mechanism comparison gap. The literature contains examples of statistical correction, temporal video processing, RGB-depth fusion, and SfM, but these mechanisms are rarely compared under one acquisition protocol with the same detector and the same tree-level truth. The appropriate experiment is not ``model A versus model B'' alone. It is an M0-to-M4 or M5 ablation in which the detector, image split, class labels, and counting truth remain fixed.

The third gap is a sensor-assumption gap. A calibrated depth camera, a point cloud, LiDAR, and monocular pseudo-depth should not be grouped under one label. Each has a different scale, failure mode, and deployment cost. An outdoor FFB benchmark should report invalid depth, registration, range, illumination, and view baseline before assigning a performance gain to the modality.

The fourth gap is a class-resolution gap. Maturity is common, but the same identity record may need multiple attributes. A future benchmark should permit at least one class attribute beyond maturity or make clear why a single attribute is the intended use case. The evaluation should distinguish class confusion from identity duplication.

The resulting research agenda is a sequence of gates. Gate~1 verifies the per-view detector and the class labels. Gate~2 builds the unique-instance tree truth and measures the M0 baseline. Gate~3 tests M1 under the fixed view protocol. Gate~4 adds temporal or appearance association. Gate~5 adds calibrated geometry or registered depth. Gate~6 tests a learned association model with a tree-disjoint split. A mechanism advances only if it lowers tree-level count error without exceeding the latency and sensor-quality budget defined before the experiment.


\section{Limitations}
\label{sec:limitations}

The search is reproducible from the captured raw exports and scoring scripts, all of which are provided as supplementary material. However, the full-text stage is intentionally targeted rather than exhaustive. Scopus was available and Web of Science was not, so database coverage is not equivalent to a two-subscription search. OpenAlex metadata and abstracts are useful for discovery but do not replace article full text. Targeted retrieval produced 44 complete evidence records from a 250-record shortlist. Consequently, the 20,035 records remain a ranked search pool, not an included-study count.

The focused evidence matrix mixes direct agricultural studies, fruit studies, transferable mechanism papers, and reviews. This is deliberate for a design-space review, but the evidence classes are not interchangeable. A result from an autonomous-driving benchmark does not establish oil-palm performance. A review that discusses ripeness does not establish cross-view identity. A dataset descriptor establishes the availability of labels and sensors, not the superiority of a model.

The review does not estimate a pooled accuracy, rank all detectors, or claim that explicit deduplication is always beneficial. Its bounded claim is that cross-view identity is a necessary design variable for the stated inventory task and that the value of each identity mechanism should be evaluated under its acquisition assumptions.


\section{Conclusion}
\label{sec:conclusion}

The target of agricultural multi-view perception is a unique, class-wise inventory, not a sum of independent detections. The search protocol produced 21,066 resolved records; 20,035 advanced after abstract screening; 250 were shortlisted for targeted review; and 44 full-text studies supplied the current evidence matrix. The direct evidence includes oil-palm detector and dataset studies, multi-angle and RGB-depth systems, fruit SfM and NeRF reconstruction, LiDAR--camera fusion, and orchard video counting. Together they support a six-mechanism design space spanning intra-view counting, statistical correction, appearance matching, temporal tracking, geometric association, and learned multi-view association.

The central contribution is a testable question rather than a fixed architecture: when does explicit identity resolution reduce tree-level duplicate and class-wise count error enough to justify its assumptions and cost? Answering that question requires tree-disjoint data, identity-linked annotations, sensor-quality reporting, and evaluation after aggregation. Oil palm is the principal case, but the design logic applies to other tree-fruit inventories whenever repeated observations must be converted into one count.


%%%%%%%%%%%%%%%%%%%%%%%%%%%%%%%%%%%%%%%%%%%%%%%%%%%%%%%%%%%%%%%%%%%%%%%%
\appendix
%%%%%%%%%%%%%%%%%%%%%%%%%%%%%%%%%%%%%%%%%%%%%%%%%%%%%%%%%%%%%%%%%%%%%%%%

\section{Focused evidence matrix}
\label{app:matrix}

Table~\ref{tab:matrix} contains one row per study included after targeted full-text review. The ``Mech.'' column uses the codes from Table~\ref{tab:mechanisms}. ``Attr.'' indicates whether the study evaluates a class attribute (e.g., maturity, size) at the instance level. ``Dedup'' indicates whether the study explicitly evaluates or performs cross-view duplicate rejection. The complete machine-readable matrix with additional fields (evidence pages, violated-assumption codes, and reviewer inferences) is available as supplementary material.

\makeatletter
\@ifclassloaded{IEEEtran}{\onecolumn}{}
\makeatother

\footnotesize
\setlength{\LTleft}{0pt}
\setlength{\LTright}{0pt}
\setlength{\tabcolsep}{2pt}
\begin{longtable}{@{}>{\raggedright\arraybackslash}p{0.17\textwidth}>{\raggedright\arraybackslash}p{0.13\textwidth}>{\raggedright\arraybackslash}p{0.07\textwidth}>{\raggedright\arraybackslash}p{0.04\textwidth}>{\raggedright\arraybackslash}p{0.04\textwidth}>{\raggedright\arraybackslash}p{0.37\textwidth}>{\raggedright\arraybackslash}p{0.06\textwidth}@{}}
\caption{Focused evidence matrix. Mechanism codes follow Table~\ref{tab:mechanisms}. Attr.\ = class attribute evaluated; Dedup = duplicate rejection evaluated or performed.}\label{tab:matrix}\\
\toprule
Study & Target, sensor & Mech. & Attr. & Dup. & Key finding and limitation & Type\\
\midrule
\endfirsthead
\multicolumn{7}{c}{\tablename\ \thetable{} continued}\\
\toprule
Study & Target, sensor & Mech. & Attr. & Dup. & Key finding and limitation & Type\\
\midrule
\endhead
\midrule
\multicolumn{7}{r}{Continued on next page}\\
\endfoot
\bottomrule
\endlastfoot

Von Hirschhausen et al.\ (2025) \cite{ledger_m_3f8f02033fe9db3a}
& Apple, calibrated stereo RGB
& M4 & -- & No
& Stereo dataset supports detection and 3D reconstruction but does not evaluate cross-view unique identity.
& Direct\\

Fusaro et al.\ (2026) \cite{ledger_r_0eda6374aaf85fae}
& Strawberry, LiDAR point cloud
& M2--4 & -- & Yes
& Learned descriptors and attention-based matching improve temporal fruit re-identification; limited to strawberry annotations.
& Direct\\

Bargoti and Underwood (2016) \cite{ledger_r_143d70deec8d2b1d}
& Apple, monocular RGB UGV
& M0 & -- & No
& CNN segmentation achieves detection F1 0.858 and row-yield $R^2$ 0.826; no persistent identity across views.
& Direct\\

Stein et al.\ (2016) \cite{ledger_r_408d13511cac4177}
& Mango, RGB + GPS/INS + LiDAR
& M3/M4 & -- & Yes
& Multi-view tracking with epipolar geometry yields near-unity calibration slope and 1.36\% over-count on 16 trees.
& Direct\\

Liu et al.\ (2018) \cite{ledger_m_9966d81433d61ed3}
& Orange/apple, monocular video + SfM
& M3/M4 & -- & Yes
& 3D reconstruction correction substantially reduces duplicate and outlier tracks; depends on illumination and camera stability.
& Direct\\

Hashim et al.\ (2018) \cite{ledger_r_97003608b250cdc5}
& Oil palm, LiDAR 905\,nm intensity
& M0 & Mat. & No
& LiDAR intensity separates three FFB ripeness levels; no classification accuracy or cross-view metric.
& Direct\\

Fawcett et al.\ (2019) \cite{ledger_m_084cfd5badbb63d3}
& Oil palm, UAV RGB SfM
& M4 & Size & Yes
& UAV canopy-height models support per-palm inventories; tree-canopy study, not fruit-level identity.
& Direct\\

Liu et al.\ (2019) \cite{ledger_r_00f82e0f404d691b}
& Mango, monocular RGB video
& M3/M4 & -- & Yes
& Monocular tracking + semantic SfM + depth filtering achieve tree-level counts; undercounting increases with occlusion.
& Direct\\

H\"ani et al.\ (2020) \cite{ledger_r_1358c0868a62c84a}
& Apple, RGB image sequences
& M3/M4 & -- & Yes
& Merged CNN counting reaches 95.6--97.8\% of harvest ground truth; inventory quality depends on detection + tracking + geometric merging composition.
& Direct\\

Alfatni et al.\ (2020) \cite{ledger_r_2f7d604b1d3fa5e4}
& Oil palm, controlled RGB CCD
& M0 & Mat. & No
& ANN histogram achieves 94\% maturity grading on conveyor; controlled setting, no multi-view inventory.
& Direct\\

Husin et al.\ (2021) \cite{ledger_r_3e45c6dc2566f46b}
& Oil palm, LiDAR point cloud
& M4 & Mat. & No
& LiDAR maps link FFB to spatial coordinates; no cross-view association or accuracy metric.
& Direct\\

Lai et al.\ (2022) \cite{lai2022yolov4palm}
& Oil palm, RGB-D RealSense D435
& M0 & Mat. & No
& YOLOv4 achieves mAP 87.9\% for ripe FFB detection at 21 FPS; no cross-frame identity.
& Direct\\

Mansour et al.\ (2022) \cite{ledger_r_19d96439ecdda18d}
& Oil palm, RGB ground images
& M0 & Mat. & No
& YOLOv5m achieves mAP$_{50{:}95}$ 0.842 on ground FFB; COVID-restricted dataset, no multi-view evaluation.
& Direct\\

Sun et al.\ (2022) \cite{ledger_r_1a70fa2aeed54ea3}
& Apple, RGB nighttime
& M0 & -- & No
& Nighttime detection AP 85.2\%; no tracking, no cross-view identity.
& Direct\\

Lu et al.\ (2022) \cite{ledger_r_2828a3937318ca86}
& Citrus, RGB handheld
& M0 & -- & No
& Improved Mask R-CNN AP$_{50}$ 95.4\% for green-fruit detection; no temporal or cross-view evaluation.
& Direct\\

Wan Nurazwin et al.\ (2022) \cite{ledger_r_3f230a23a2a956c8}
& Pineapple, UAV top-view RGB
& M0 & -- & No
& ANN counting achieves 94.4\% on 20 sample images; no cross-view deduplication.
& Direct\\

Yang et al.\ (2022) \cite{ledger_r_4b3d8d423dbd5d0e}
& Vehicle, RGB multi-camera
& M2 & -- & N/A
& Weakly supervised cross-view embeddings reduce viewpoint variation; transferable mechanism, not agricultural data.
& Transfer\\

Zhang et al.\ (2022a) \cite{ledger_r_c79d61a639a09929}
& Apple, UAV + ground SfM
& M4 & -- & No
& Complementary UAV and ground point clouds improve coverage; no persistent fruit identity or deduplication.
& Direct\\

Zhang et al.\ (2022b) \cite{ledger_r_ee550a3e3a6fb318}
& Citrus, RGB video rover
& M3 & -- & Yes
& Counting-region strategy with tracking achieves MAE 0.081; track IDs are local to each sequence.
& Direct\\

Chen et al.\ (2022) \cite{chen2022ffbdet}
& Oil palm, RGB online images
& M0 & -- & No
& YOLOv4 mAP 98.5\% on small dataset; Kaggle/Google-sourced images, no field or multi-view validation.
& Direct\\

Phimpisan and Chungchoo (2023) \cite{ledger_r_1c81e223679e4bfe}
& Oil palm, UV fluorescence
& M0 & Mat. & No
& Fluorescence classifies FFB ripeness rapidly; controlled dark-chamber setup, no field multi-view use.
& Direct\\

Gené-Molá et al.\ (2023) \cite{ledger_r_2d972b94554c6cd1}
& Apple, Kinect Azure RGB-D / video
& M3 & -- & Yes
& ByteTrack HOTA 0.689, yield MAPE 7.47\%; appearance ReID did not help visually similar apples.
& Direct\\

Shiddiq et al.\ (2023) \cite{shiddiq2023palmcount}
& Oil palm, RGB video conveyor
& M3 & -- & No
& Conveyor counting mean accuracy 79.4\%; postharvest setup, sensitive to illumination and orientation.
& Direct\\

Rapado-Rincón et al.\ (2023) \cite{ledger_r_700d34dae0f20d78}
& Tomato, RGB + LiDAR robot arm
& M3/M4 & -- & Yes
& Multi-view perception recovers occluded fruit; HOTA up to 71.5\%, MAPE 12.2\%. Count error alone hides ID switches.
& Direct\\

Abeyrathna et al.\ (2023) \cite{ledger_r_9e1316b934e66409}
& Apple, RGB-D RealSense + tractor
& M3 & -- & Yes
& Deep SORT with depth prevents repeated counting; artificial orchard with 30 apples on 10 trees.
& Direct\\

Hu et al.\ (2023) \cite{ledger_r_d594d85dc382f9ff}
& Apple, RGB video
& M3 & -- & Yes
& SURF + Kalman cascade tracking reduces repeated counts vs SORT; three test videos, no full identity metric.
& Direct\\

Lai et al.\ (2023) \cite{lai2023palmslr}
& Oil palm, multi-sensor review
& M0 & Mat. & N/A
& Review covers ripeness and detection; does not formulate cross-view identity or unique inventory as a problem.
& Direct\\

James et al.\ (2023) \cite{ledger_r_fd5751e2ca2f7496}
& Citrus, RGB ground images
& M0 & -- & No
& Public citrus benchmark with 44,233 boxes; front/back aggregation does not establish unique fruit identity.
& Direct\\

Lin et al.\ (2024) \cite{ledger_r_1c69e46dac45d45d}
& Citrus, RGB ground images
& M0 & -- & No
& AG-YOLO mAP$_{50}$ 83.2\% under severe occlusion; single-view benchmark, no cross-view identity.
& Direct\\

Ambrus et al.\ (2024) \cite{ledger_r_3e129b7285059b95}
& Tomato, DSLR + robot RGB + 3D scan
& M0/M4 & Mat. & No
& Robot/DSLR + 3D shape supports count and weight estimation; one-sided views cause missed fruit.
& Direct\\

V\'elez et al.\ (2024) \cite{ledger_r_659a7da54b8a43e6}
& Grape, smartphone + UAV + RTK
& M4 & -- & N/A
& Geospatial anchors align multi-source data at plant level; no individual fruit identity.
& Transfer\\

Wu et al.\ (2024) \cite{ledger_r_92ac7cb4b797197d}
& Apple, RGB orchard + synthetic fog
& M0 & -- & No
& Lightweight detector mAP$_{50}$ 94.3\% across weather conditions; no tracking or cross-view identity.
& Direct\\

Vidhya and Priya (2024) \cite{ledger_r_e8d88a31640ee300}
& Mango, RGB mobile camera
& M0 & -- & No
& YOLOv7 mAP$_{50}$ 99.5\% reported on training/validation; no separate test or cross-view metrics.
& Direct\\

Hobart et al.\ (2025) \cite{ledger_m_66042c905cc191de}
& Apple, UAV RGB SfM + RTK
& M4 & -- & No
& SfM recovers ripe-apple 3D positions for yield-mass mapping; no persistent fruit identity.
& Direct\\

Montalb\'an Faet et al.\ (2025) \cite{ledger_m_949f38a7336860d8}
& Citrus, UAV RGB + multispectral
& M4 & -- & No
& Cross-spectral annotations with date-based split; data article without detector benchmark.
& Direct\\

Dutta et al.\ (2025) \cite{ledger_m_dcfc5d3cb912146d}
& Apple, UAV + ESP32-CAM edge
& M0 & -- & No
& Low-cost edge pipeline for detection; temporal consistency identified as future work.
& Direct\\

Sabillah et al.\ (2025) \cite{ledger_r_134385f2faed38a6}
& Oil palm, RGB smartphone + UAV
& M0 & -- & No
& YOLOv8 mAP 99.0\%; small dataset, no independent test split or multi-view evaluation.
& Direct\\

Goh et al.\ (2025) \cite{goh2025palmdataset}
& Oil palm, RGB-D RealSense + point cloud
& M4 & Mat. & No
& 400 RGB-depth-point-cloud scenes with calibration; data article, no detector benchmark or cross-session identity.
& Direct\\

Timprae et al.\ (2025) \cite{ledger_r_4d09eb83eea41b2f}
& Tomato, RGB handheld + SfM/NeRF
& M4 & Mat. & No
& Instance segmentation + SfM + DBSCAN + ellipsoid fitting for phenotyping; manual intervention still required.
& Direct\\

Hamid et al.\ (2025) \cite{ledger_r_504a1d386928344f}
& Oil palm, multi-sensor review
& M0 & Mat. & N/A
& Mini-review on ripeness; narrative format, no reproducible search or cross-view analysis.
& Direct\\

Shiddiq et al.\ (2025) \cite{ledger_r_722bfce242f37ecb}
& Oil palm, RGB CMOS conveyor
& M0 & Mat. & No
& Front/back ImageJ counts are lower than manual due to overlapping fruitlets; no cross-view registration.
& Direct\\

Lee et al.\ (2026) \cite{ledger_r_d4270ea8ab11f268}
& Pear, RGB + LiDAR + SLAM
& M3/M4 & -- & Yes
& 96.2\% counting accuracy with LiDAR--camera fusion; 53 double-counts remain from threshold and SLAM drift.
& Direct\\

Meyer et al.\ (2024) \cite{ledger_m_1210e3ea00486b68}
& Fruit (synthetic + apple), RGB multiview NeRF
& M4 & -- & Yes
& Semantic NeRF clusters fruit in 3D, addressing duplicate counting; sensitive to sparse views and hyperparameters.
& Direct\\

Wang et al.\ (2025) \cite{ledger_r_b2e0e9aabad17a97}
& Tomato, RGB-D + ORB-SLAM3
& M4 & Size & No
& RGB-D SLAM projects segmentation into 3D for per-plant count and volume; leaf occlusion and depth errors cause missed fruit.
& Direct\\

\end{longtable}

\makeatletter
\@ifclassloaded{IEEEtran}{\twocolumn}{}
\makeatother


\section{Exact search strings}
\label{app:queries}

The following blocks reproduce the final Scopus strings. All seven use \texttt{PUBYEAR > 2014 AND PUBYEAR < 2027} and a subject-area clause covering Agriculture, Computer Science, Engineering, and Multidisciplinary (Earth Sciences and Environmental Science are added in Q3 and Q6). The line breaks are for typesetting only.

\subsection*{Q1: Multi-view fruit inventory}
\begin{verbatim}
TITLE-ABS-KEY(("fruit" OR "fresh fruit
bunch" OR "FFB" OR "oil palm" OR "apple"
OR "citrus" OR "mango" OR "grape" OR
"berry" OR "bunch" OR "cluster" OR "crop")
AND ("count*" OR "yield estimation" OR
"load estimation" OR "fruit load" OR
"inventory" OR "enumeration")
AND ("multi-view" OR "multiple view*" OR
"multi-camera" OR "video" OR "cross-view"
OR "structure from motion" OR "SfM" OR
"track*" OR "3D reconstruction" OR
"image sequence"))
AND PUBYEAR > 2014 AND PUBYEAR < 2027
AND (SUBJAREA(AGRI) OR SUBJAREA(COMP)
OR SUBJAREA(ENGI) OR SUBJAREA(MULT))
\end{verbatim}

\subsection*{Q2: Cross-view identity and duplicates}
\begin{verbatim}
TITLE-ABS-KEY(("multi-view" OR "cross-view"
OR "multi-camera" OR "multiple cameras" OR
"overlapping view*" OR "multi-sensor")
AND ("re-identification" OR
"data association" OR "identity" OR
"duplicate*" OR "deduplicat*" OR
"double counting" OR "correspondence" OR
"instance matching" OR "track*" OR
"association")
AND ("object detection" OR
"instance segmentation" OR "counting" OR
"instance" OR "pedestrian" OR "vehicle"))
AND PUBYEAR > 2014 AND PUBYEAR < 2027
AND (SUBJAREA(COMP) OR SUBJAREA(ENGI)
OR SUBJAREA(AGRI) OR SUBJAREA(MULT))
\end{verbatim}

\subsection*{Q3: Plant geometry and multimodal sensing}
\begin{verbatim}
TITLE-ABS-KEY(("RGB-D" OR "depth camera"
OR "stereo vision" OR "LiDAR" OR
"point cloud" OR "monocular depth" OR
"photogrammetry" OR
"structure from motion")
AND ("fruit" OR "orchard" OR "vineyard"
OR "canopy" OR "plant" OR "tree crop" OR
"plantation" OR "oil palm")
AND ("detection" OR "segmentation" OR
"localization" OR "counting" OR
"phenotyping" OR "harvesting"))
AND PUBYEAR > 2014 AND PUBYEAR < 2027
AND (SUBJAREA(AGRI) OR SUBJAREA(COMP)
OR SUBJAREA(ENGI) OR SUBJAREA(EART)
OR SUBJAREA(ENVI) OR SUBJAREA(MULT))
\end{verbatim}

\subsection*{Q4: Per-instance attributes}
\begin{verbatim}
TITLE-ABS-KEY(("fruit" OR "produce" OR
"crop" OR "bunch")
AND ("maturity" OR "ripeness" OR "grading"
OR "grade" OR "quality class*" OR
"size class*" OR "defect" OR "disease" OR
"cultivar" OR "variety")
AND ("instance segmentation" OR
"object detection" OR "per-instance" OR
"individual fruit" OR "bounding box" OR
"per-object"))
AND PUBYEAR > 2014 AND PUBYEAR < 2027
AND (SUBJAREA(AGRI) OR SUBJAREA(COMP)
OR SUBJAREA(ENGI) OR SUBJAREA(MULT))
\end{verbatim}

\subsection*{Q5: Previous reviews}
\begin{verbatim}
TITLE(("review" OR "survey" OR
"systematic review" OR "scoping review"))
AND TITLE-ABS-KEY(("fruit detection" OR
"fruit counting" OR "yield estimation" OR
"oil palm" OR "orchard" OR
"crop monitoring" OR
"agricultural robot*" OR
"precision agriculture")
AND ("deep learning" OR "computer vision"
OR "object detection" OR
"machine vision"))
AND PUBYEAR > 2014 AND PUBYEAR < 2027
AND (SUBJAREA(AGRI) OR SUBJAREA(COMP)
OR SUBJAREA(ENGI) OR SUBJAREA(MULT))
\end{verbatim}

\subsection*{Q6: Oil-palm vision and machine learning}
\begin{verbatim}
TITLE-ABS-KEY(("oil palm" OR
"elaeis guineensis" OR
"fresh fruit bunch" OR "FFB")
AND ("deep learning" OR "machine learning"
OR "convolutional" OR "neural network" OR
"computer vision" OR "image processing" OR
"YOLO" OR "object detection" OR
"remote sensing"))
AND PUBYEAR > 2014 AND PUBYEAR < 2027
AND (SUBJAREA(AGRI) OR SUBJAREA(COMP)
OR SUBJAREA(ENGI) OR SUBJAREA(EART)
OR SUBJAREA(ENVI) OR SUBJAREA(MULT))
\end{verbatim}

\subsection*{Q7: Single-view fruit load}
\begin{verbatim}
TITLE-ABS-KEY(("fruit" OR "apple" OR
"citrus" OR "mango" OR "grape" OR "berry"
OR "bunch" OR "oil palm" OR "orchard" OR
"vineyard")
AND ("counting" OR "count" OR
"yield estimation" OR "load estimation" OR
"fruit load" OR "crop load")
AND ("deep learning" OR "convolutional" OR
"object detection" OR "YOLO" OR
"instance segmentation" OR
"machine vision"))
AND PUBYEAR > 2014 AND PUBYEAR < 2027
AND (SUBJAREA(AGRI) OR SUBJAREA(COMP)
OR SUBJAREA(ENGI) OR SUBJAREA(MULT))
\end{verbatim}
